%
% graph.tex -- Template für standalone TIKZ Bilder
%
% (c) 2019 Prof Dr Andreas Müller, Hochschule Rapperswil
%
\documentclass[tikz,12pt]{standalone}
\usepackage{amsmath}
\usepackage{times}
\usepackage{txfonts}
\usepackage{pgfplots}
\usepackage{csvsimple}
\usetikzlibrary{arrows,intersections,math,calc}
\definecolor{darkred}{rgb}{0.8,0,0}
\definecolor{darkgreen}{rgb}{0,0.6,0}
\def\Cminus{0.5348}
\def\Cplus{-1.2279}
\def\a{0.785398}
\begin{document}
\def\skala{4}
\begin{tikzpicture}[>=latex,thick,scale=\skala,
declare function = {
w(\X) = 180 * \X / 3.1415;
fplus(\X,\C)  = -ln(cos(w(\X)))+0.5*ln((1+sin(w(\X)))/(1-sin(w(\X))))+(\C);
fminus(\X,\C) = -ln(cos(w(\X)))-0.5*ln((1+sin(w(\X)))/(1-sin(w(\X))))+(\C);
}]

\fill[color=gray!20] (0,-1.5) rectangle (1,1);
\draw[color=gray,line width=0.2pt] (1,-1.5) -- (1,1);
\node[color=white] at (0.6,-1.1) [scale=6] {$\Omega$};
\node[color=white] at (0.6,0.6) [scale=6] {$\Omega$};

\fill[color=blue!10]
	plot[domain=0:\a,samples=20] (\x,{fplus(\x,\Cplus)})
	--
	plot[domain=0:\a,samples=20] ({\a-\x},{fminus((\a-\x),\Cminus)})
	--
	cycle;
\draw[color=blue]
	plot[domain=0:\a,samples=20] (\x,{fplus(\x,\Cplus)})
	--
	plot[domain=0:\a,samples=20] ({\a-\x},{fminus((\a-\x),\Cminus)});

\begin{scope}
\clip (0,-1.5) rectangle (1.1,1.0);
\foreach \c in {-3.6,-3.4,...,1.0}{
	\draw[color=darkgreen]
		plot[domain=0:1.1,samples=100] ({\x},{fplus(\x,\c)});
}
\foreach \c in {-1.4,-1.2,...,1.6}{
	\draw[color=darkred]
		plot[domain=0:1.1,samples=100] ({\x},{fminus(\x,\c)});
}
\end{scope}

\draw[->] (-0.02,0) -- (1.4,0) coordinate[label={$x$}];
\draw[->] (0,-1.52) -- (0,1.06) coordinate[label={left:$y$}];
\node at (-0.02,0) [left] {$0$};

\fill[color=blue] (\a,0) circle[radius=0.0125];
\node[color=blue] at (\a,0) [above right] {$P$};

\fill[color=blue] (0,\Cplus) circle[radius=0.0125];
\node[color=blue] at (0,\Cplus) [left] {$C_+$};
\fill[color=blue] (0,\Cminus) circle[radius=0.0125];
\node[color=blue] at (0,\Cminus) [left] {$C_-$};

\end{tikzpicture}
\end{document}

